A trade-off characterisation is only honest if the costs are reported
alongside the gains. We document two small but consistent negatives.

\subsection{Out-of-distribution perplexity}
\label{sec:results-ood}

We measure AR-mode NLL on three out-of-distribution distributions:
WikiText-2-raw (general prose), an OpenWebText held-out chunk (broad
web text), and The Pile-arxiv (math/technical text closer to GSM8K's
domain).

\begin{table}[h]
\centering
\begin{tabular}{lrrrrr}
\toprule
Scale & comp.\ wikitext & ar.\ wikitext & comp.\ owt & ar.\ owt & $\Delta$ wikitext / owt \\
\midrule
$60$M & 8.90 & 8.78 & 8.31 & 8.30 & $+0.12 / +0.01$ \\
$120$M & 8.99 & 9.01 & 8.35 & 8.36 & $-0.02 / -0.01$ \\
$200$M & 9.09 & 9.00 & 8.37 & 8.35 & $+0.09 / +0.02$ \\
$250$M & 9.21 & 9.05 & 8.35 & 8.28 & $+0.16 / +0.07$ \\
$300$M & 9.36 & 9.18 & 8.58 & 8.38 & $+0.18 / +0.20$ \\
\bottomrule
\end{tabular}
\caption{AR-mode NLL on OOD chunks (WikiText-2-raw and OpenWebText
held-out). Composite is uniformly $+0.1$--$0.2$ NLL worse on prose,
with the gap widening at larger scale.}
\label{tab:ood}
\end{table}

\textbf{The composite trade does not transfer to OOD prose.} The
regularising effect that helps the diffusion head and the
data-constrained AR head does not produce a more generalising AR
representation on non-math English text. The penalty is small in
absolute terms ($\le 0.2$ NLL) but consistent and growing with
scale.

\subsection{Calibration}
\label{sec:results-ece}

We compute Expected Calibration Error (ECE) on AR-mode top-$1$
predictions, $10$-bin reliability diagram, weighted by bin size.

\begin{table}[h]
\centering
\begin{tabular}{lrrrr}
\toprule
Scale & comp.\ ECE & ar.\ ECE & comp.\ top-$1$ & ar.\ top-$1$ \\
\midrule
$60$M & 0.029 & 0.024 & 0.59 & 0.62 \\
$120$M & 0.028 & 0.027 & 0.60 & 0.64 \\
$200$M & 0.026 & 0.027 & 0.48 & 0.52 \\
$250$M & 0.024 & 0.020 & 0.44 & 0.49 \\
$300$M & 0.022 & 0.013 & 0.38 & 0.46 \\
\bottomrule
\end{tabular}
\caption{Expected Calibration Error (lower is better) and top-$1$
accuracy on AR-mode predictions on GSM8K held-out answer tokens.
Composite is less confident and less accurate at every scale, with
ECE worse at $4$ of $5$ scales.}
\label{tab:ece}
\end{table}

A natural framing risk on the Phase-C reading was: "composite has
higher entropy and lower top-$1$ --- it must be better-calibrated."
The proper ECE calculation falsifies this: composite is \emph{less
confident and less accurate}, which produces \emph{worse} calibration
(higher ECE) at most scales.

\subsection{Combined reading}

Both negatives are small in absolute terms. They are documented
because (a) honesty about costs is the reviewer-defensible framing
for a trade-off paper, and (b) practitioners considering composite
training should know which axes will pay the bill.

The full claim ledger for this study:

\begin{table}[h]
\centering
\begin{tabular}{lll}
\toprule
Claim & Confidence & Evidence \\
\midrule
Joint training has real diff-axis advantage & High &
\textsection\ref{sec:results-compute} \\
Data-regime crossover at $\sim 1$k problems & High &
\textsection\ref{sec:results-curve} \\
Mode-switching gives small downstream lift & Medium &
\textsection\ref{sec:results-modeswitch}, $n=3 \to 8$ \\
Composite is better at OOD prose & \textbf{NO} & \textsection\ref{sec:results-ood} \\
Composite is better-calibrated & \textbf{NO} & \textsection\ref{sec:results-ece} \\
\bottomrule
\end{tabular}
\caption{The claim ledger. We make three trade-off claims and document
two negatives.}
\label{tab:claim-ledger}
\end{table}
